% Generated draft: requires mathematical and source review.
\documentclass[11pt]{article}
\usepackage[margin=25mm]{geometry}
\usepackage{fontspec}
\setmainfont{Latin Modern Roman}
\usepackage{amsmath,amssymb,mathtools,unicode-math}
\setmathfont{Latin Modern Math}
\usepackage{enumitem,needspace,xurl,hyperref}
\setlength{\parindent}{0pt}
\setlength{\parskip}{5pt}
\setlength{\emergencystretch}{3em}
\newcommand{\N}{\mathbb{N}}
\newcommand{\Z}{\mathbb{Z}}
\newcommand{\Q}{\mathbb{Q}}
\newcommand{\R}{\mathbb{R}}
\newcommand{\C}{\mathbb{C}}
\newcommand{\F}{\mathbb{F}}
\newcommand{\A}{\mathbb{A}}
\newcommand{\PP}{\mathbb{P}}
\newcommand{\E}{\mathbb{E}}
\newcommand{\eps}{\varepsilon}
\newcommand{\dd}{\,\mathrm{d}}
\begin{document}
\section*{UnsolvedMath ID: 20000557}
\textbf{Source problem code:} catalogue entry
\subsection{Definitions and mathematical statement}
For \(1\leq p\leq\infty\), let
\[
 \|A\|_{p\to p}=\sup_{x\neq0}\frac{\|Ax\|_{\ell_p^n}}{\|x\|_{\ell_p^n}}
\]
be the induced operator norm, not a Schatten norm. For \(\sigma\subset\{1,\ldots,n\}\), let \(P_\sigma\) be coordinate projection. If \(\mathcal P=(\sigma_1,\ldots,\sigma_k)\) is a partition, with empty blocks allowed, put \(A_{\mathcal P}=\sum_jP_{\sigma_j}AP_{\sigma_j}\). Then \(\|A_{\mathcal P}\|_{p\to p}=\max_j\|P_{\sigma_j}AP_{\sigma_j}\|_{p\to p}\).

Bill Johnson's source question asks, for every fixed \(1\leq p<\infty\), whether
\[
 \forall\epsilon>0\ \exists k=k(\epsilon,p)<\infty\ \forall n\geq1\ \forall A\in M_n(\C),\quad
 \left[ A_{ii}=0,\ \|A\|_{p\to p}=1\right]
 \Longrightarrow\exists\mathcal P\text{ with at most }k\text{ blocks}:\
 \|A_{\mathcal P}\|_{p\to p}\leq\epsilon.
\]
The source does not specify real versus complex scalars; the displayed convention uses complex matrices. Let \(R_p(\epsilon)\) be the least such number, or \(\infty\) if none exists, extending the definition to \(p=\infty\) for the supplied duality statement. For \(0<\epsilon<1\), the aid reports
\[
 R_1(\epsilon)\leq\lceil2/\epsilon\rceil,\qquad
 R_2(\epsilon)=\Theta(\epsilon^{-2})\quad(\epsilon\downarrow0),\qquad
 R_p(\epsilon)=R_{p'}(\epsilon),\quad \frac1p+\frac1{p'}=1,
\]
with \(1'=\infty\), \(\infty'=1\). Its literature search located no general settlement for \(1<p<\infty\), \(p\neq2\); this is a supplied search qualification rather than a verified current nonexistence of a theorem.

The separate simultaneous-paving theorem concerns a real symmetric conference matrix \(C\) of order \(n\geq2\): \(C_{ii}=0\), \(C_{ij}\in\{-1,1\}\) for \(i\neq j\), and \(C^2=(n-1)I_n\). The aid states that, for every \(0<\epsilon<1\), there is a single partition \(\mathcal P\) into at most \(\lceil54\epsilon^{-2}\rceil\) blocks such that
\[
 \|C_{\mathcal P}\|_{p\to p}\leq\epsilon\|C\|_{p\to p}
 \qquad\text{simultaneously for every }1\leq p\leq\infty.
\]
It obtains this from one Ravichandran--Srivastava multi-paving applied to the three Hermitian contractions
\[
 C/\sqrt{n-1},\qquad -C/\sqrt{n-1},\qquad (J-I_n)/(n-1),
\]
where \(J\) is the all-ones matrix. The statement does not assert such simultaneous partitions for arbitrary matrices. The aid further reports that the infinite Paley conference family forces every fixed-\(p\) uniform paving number to satisfy
\[
 R_p(\epsilon)\geq\epsilon^{-\min(p,p')}\qquad(0<\epsilon<1),
\]
including the endpoint convention above. This is a necessary bound for a uniform number valid over the family, not a claim that every individual finite matrix requires that many blocks.

Finally the aid reports an explicit rank-one direct-sum family which is exactly two-pavable while naive endpoint interpolation has unbounded dimension-dependent defect. Its matrices and the precise defect quantity are not supplied here, so no more detailed interpolation theorem is inferred. The original comments report \(\mathrm{Pav}_1\) as long known, tentatively attributed to Bourgain, compare it informally with nonnegative-entry \(\mathrm{Pav}_2\), and recall Bourgain--Tzafriri restricted invertibility with optimal estimates only at \(p=2\) in their original work.
\subsection{Short English statement}
The question asks for a dimension-independent coordinate paving bound in induced \(\ell_p\) operator norm. The supplied report gives the endpoint results and duality, while describing the general intermediate case as unlocated in its search. A separate theorem gives one partition that paves each real symmetric conference matrix in every \(\ell_p\) norm, with a Paley-family lower bound.
\subsection{Sources}
\begin{itemize}
\item Local registry: permanent UnsolvedMath ID \(20000557\), exact public code \texttt{catalogue entry}; \texttt{aim-analysis-notes.json}, \texttt{notes[88]}.
\item AIM workshop \emph{Beyond Kadison--Singer: paving and consequences}, item b16, attributed in the supplied verification field to Bill Johnson: \url{https://docs.google.com/document/d/1XRw0WyG5PshqboCp99tovCRD1GCK0jXZiNLH_nw618c/edit}.
\item Kenneth Berman, Herbert Halpern, Victor Kaftal, Gary Weiss, ``Matrix norm inequalities and the relative Dixmier property,'' \emph{Integral Equations Operator Theory} 11 (1988), no. 1, 28--48. The source relates this to nonnegative-matrix paving.
\item The original text tentatively attributes \(\mathrm{Pav}_1\) to Bourgain using ``IIRC,'' and mentions the original Bourgain--Tzafriri restricted-invertibility paper without a complete citation.
\item Supplied research summary, ``Induced ell-p paving endpoints and simultaneous all-p conference paving,'' attempt 2, cites Ravichandran--Srivastava multi-paving and the Paley family without full bibliographic entries. The original literature field is empty.
\end{itemize}
\end{document}
